\section{CertGraph Static Architecture and Bilinear Gating}
\label{sec:certgraph_arch}

This section details the static detection layer of CertGraph (Figure~\ref{fig:certgraph_arch}), analyzes the theoretical failure of standard message passing on access control graphs, and introduces a bilinear compatibility formulation whose learned output must be distinguished from exact authorization reachability.

\subsection{Heterogeneous Graph Attention Network (HetGAT)}
Following graph-learning formulations~\cite{scarselli2008graph,wu2021survey}, CertGraph employs a multi-relational graph attention network. Let $\mathbf{x}_v \in \mathbb{R}^{d_{\phi(v)}}$ denote the initial feature vector of entity $v \in \mathcal{V}$. Vertices are first projected into a shared hidden space of dimension $d$:
\begin{equation}
\fitmath{\mathbf{h}_v^{(0)} = \text{ReLU} \left( \mathbf{W}_{\phi(v)}^{(0)} \mathbf{x}_v + \mathbf{b}_{\phi(v)}^{(0)} \right)}
\end{equation}
where $\mathbf{W}_{\phi(v)}^{(0)} \in \mathbb{R}^{d \times d_{\phi(v)}}$ is a type-specific linear projection matrix.

At each layer $l \in \{1, \dots, L\}$, message passing aggregates relational representations across incoming edges $\mathcal{N}_r(v) = \{u \in \mathcal{V} \mid (u, v) \in \mathcal{E}, \psi(u, v) = r\}$:
\begin{equation}
\fitmath{\mathbf{h}_v^{(l)} = \sigma \left( \sum_{r \in \mathcal{R}_{\text{in}}(\phi(v))} \sum_{u \in \mathcal{N}_r(v)} \alpha_{vu}^{(l)} \mathbf{W}_r^{(l)} \mathbf{h}_u^{(l-1)} \right)}
\end{equation}
where $\mathbf{W}_r^{(l)} \in \mathbb{R}^{d \times d}$ is the relation-specific weight transformation for relation $r \in \mathcal{T}_E$, and $\sigma(\cdot)$ denotes the ELU activation function. The relational attention coefficient $\alpha_{vu}^{(l)}$ assigns dynamic importance to edge $(u, v)$ under relation $r$:
\begin{equation}
\fitmath{\begin{gathered}
\alpha_{vu}^{(l)}=\frac{\exp(z_{vu}^{(l)})}{\sum_{k\in\mathcal N_r(v)}\exp(z_{vk}^{(l)})},\\
z_{vu}^{(l)}=\operatorname{LeakyReLU}\!\left(\mathbf a_r^\top
\begin{bmatrix}\mathbf W_r^{(l)}\mathbf h_v^{(l-1)}\\ \mathbf W_r^{(l)}\mathbf h_u^{(l-1)}\end{bmatrix}\right).
\end{gathered}}
\end{equation}
where $\mathbf{a}_r \in \mathbb{R}^{2d}$ is the parameterized attention vector for relation $r$, and $\|$ denotes vector concatenation.

\begin{figure*}[ht]
\centering
\includegraphics[width=\textwidth]{figures/certgraph_learning_architecture.pdf}
\caption{CertGraph learning architecture and the literal bilinear mechanism. (A) Typed graph features enter (B) two relational attention layers with four heads per layer and gated self-skips, followed by (C) template readout and binary classification in the matched resolved-permission ESC1 study. (D) Source and target states determine a scalar bilinear gate that scales the transformed neighbor message before attention-weighted aggregation. Exact authorization is checked separately; the learned gate is not a Boolean reachability test. The input topology is illustrative.}
\label{fig:certgraph_arch}
\end{figure*}

\subsection{Theoretical Analysis: Representation Collapse and Gradient Bounds}
One limitation in pure neighborhood aggregation on directed multigraphs is the erasure of intrinsic attributes for leaf and intermediate nodes.

\begin{theorem}[Intrinsic Attribute Preservation and Gradient Lower Bounds]
\label{thm:theorem1}
Let $\mathcal{G} = (\mathcal{V}, \mathcal{E}, \phi, \psi)$ be a directed heterogeneous multigraph. Consider an $L$-layer pure relational aggregation model with no target-feature dependence in its aggregation coefficients, no biases, and $\sigma(\mathbf{0})=\mathbf{0}$. For the residual bound, differentiate with respect to the common hidden input $\mathbf h_v^{(0)}$ and assume square compatible skip maps and differentiability. No lower bound for the raw type-projected features is asserted. These assumptions restrict the theorem; attention coefficients that depend on the target features can transmit those features even without a self-loop.
\begin{enumerate}
    \item \textbf{Intrinsic Feature Erasure:} If $\mathcal{G}$ contains no explicit self-loops on template nodes ($v \notin \mathcal{N}_r(v)$ for all $r$), the first-layer representation $\mathbf{h}_v^{(1)}$ is conditionally independent of initial features $\mathbf{x}_v$:
    \begin{equation}
\fitmath{\frac{\partial \mathbf{h}_v^{(1)}}{\partial \mathbf{x}_v} = \mathbf{0} \in \mathbb{R}^{d \times d_{\phi(v)}}}
\end{equation}
    Furthermore, for any unprivileged source account $u$ with in-degree $d_{\text{in}}(u) = 0$, $\mathbf{h}_u^{(l)} = \mathbf{0}$ for all $l \ge 1$.
    \item \textbf{Guaranteed Gradient Lower Bound via Residual Skips:} When parameterized residual skip connections are introduced ($\tilde{\mathbf{h}}_v^{(l)} = \mathbf{h}_v^{(l)} + \mathbf{W}_{\text{skip}}^{(l)} \tilde{\mathbf{h}}_v^{(l-1)}$) with $\sigma_{\min}(\mathbf{W}_{\text{skip}}^{(l)}) > 0$, the Jacobian of the representation satisfies:
    \begin{equation}
\fitmath{\sigma_{\min} \left( \frac{\partial \tilde{\mathbf{h}}_v^{(L)}}{\partial \mathbf{h}_v^{(0)}} \right) \ge \prod_{k=1}^L \sigma_{\min}(\mathbf{W}_{\text{skip}}^{(k)}) - \|\mathcal{J}_{\text{graph}}(v)\|}
\end{equation}
    where $\mathcal{J}_{\text{graph}}(v)$ is the full derivative of all non-skip contributions, including their dependence through earlier residual states. Its norm must be measured or bounded for the particular model; a single path's weight product does not bound an arbitrary graph's aggregate derivative. When $\prod_{k=1}^L \sigma_{\min}(\mathbf{W}_{\text{skip}}^{(k)}) > \|\mathcal{J}_{\text{graph}}(v)\|$, initial feature sensitivity is strictly preserved:
    \begin{equation}
\fitmath{\left\| \frac{\partial \tilde{\mathbf{h}}_v^{(L)}}{\partial \mathbf{h}_v^{(0)}} \right\| > 0}
\end{equation}
\end{enumerate}
\end{theorem}

\begin{proof}
We establish the two claims sequentially:

\paragraph{Part 1: Intrinsic Feature Erasure under Pure Neighborhood Aggregation.}
Under pure relational convolution at layer $l=1$, the hidden state $\mathbf{h}_v^{(1)}$ is computed as $\sigma \left( \sum_{r \in \mathcal{R}_{\text{in}}(\phi(v))} \sum_{u \in \mathcal{N}_r(v)} \alpha_{vu}^{(1)} \mathbf{W}_r^{(1)} \mathbf{h}_u^{(0)} \right)$. In an Active Directory multigraph, template nodes possess no explicit self-loops ($v \notin \mathcal{N}_r(v)$ for all $r$). Differentiating $\mathbf{h}_v^{(1)}$ with respect to the initial template feature vector $\mathbf{x}_v$ through its type-specific initial projection:
\begin{equation}
\fitmath{\begin{gathered}\frac{\partial \mathbf{h}_v^{(1)}}{\partial \mathbf{x}_v} \\ = \sigma'(\cdot) \sum_{r \in \mathcal{R}_{\text{in}}(\phi(v))} \sum_{u \in \mathcal{N}_r(v)} \alpha_{vu}^{(1)} \mathbf{W}_r^{(1)} \frac{\partial \mathbf{h}_u^{(0)}}{\partial \mathbf{x}_v}\end{gathered}}
\end{equation}
Because distinct nodes possess independent initial feature vectors ($\partial \mathbf{h}_u^{(0)} / \partial \mathbf{x}_v = \mathbf{0}$ for all $u \neq v$), every term in the summation vanishes identically, yielding $\partial \mathbf{h}_v^{(1)} / \partial \mathbf{x}_v = \mathbf{0}$. Thus, the layer-1 embedding is conditionally independent of $\mathbf{x}_v$, discarding all intrinsic template configuration flags (e.g., \texttt{CT\_FLAG\_ENROLLEE\_SUPPLIES\_SUBJECT}). Moreover, for any source node $u$ with in-degree $d_{\text{in}}(u) = 0$, $\mathcal{N}_r(u) = \emptyset$ for all $r$, so $\mathbf{h}_u^{(l)} = \mathbf{0}$ for all $l \ge 1$ by empty summation and $\sigma(\mathbf{0}) = \mathbf{0}$.

\paragraph{Part 2: Gradient Lower Bounds via Parameterized Residual Skips.}
With parameterized residual skip connections $\tilde{\mathbf{h}}_v^{(l)} = \mathbf{h}_v^{(l)} + \mathbf{W}_{\text{skip}}^{(l)} \tilde{\mathbf{h}}_v^{(l-1)}$ and $\tilde{\mathbf{h}}_v^{(0)} = \mathbf{h}_v^{(0)}$, unrolling the recurrence across $L$ layers yields:
\begin{equation}
\fitmath{\begin{gathered}\tilde{\mathbf{h}}_v^{(L)} = \left( \prod_{k=1}^L \mathbf{W}_{\text{skip}}^{(k)} \right) \mathbf{h}_v^{(0)} \\ + \sum_{l=1}^L \left( \prod_{k=l+1}^L \mathbf{W}_{\text{skip}}^{(k)} \right) \mathbf{h}_v^{(l)}\end{gathered}}
\end{equation}
Differentiating $\tilde{\mathbf{h}}_v^{(L)}$ with respect to $\mathbf{h}_v^{(0)}$:
\begin{equation}
\fitmath{\begin{gathered}\frac{\partial \tilde{\mathbf{h}}_v^{(L)}}{\partial \mathbf{h}_v^{(0)}} = \prod_{k=1}^L \mathbf{W}_{\text{skip}}^{(k)} + \mathcal{J}_{\text{graph}}(v),\\ \text{where}\quad \mathcal{J}_{\text{graph}}(v) = \sum_{l=1}^L \left( \prod_{k=l+1}^L \mathbf{W}_{\text{skip}}^{(k)} \right) \frac{\partial \mathbf{h}_v^{(l)}}{\partial \mathbf{h}_v^{(0)}}\end{gathered}}
\end{equation}
By Weyl's singular value perturbation inequality and the sub-multiplicative property of singular values:
\begin{equation}
\fitmath{\begin{gathered}\sigma_{\min} \left( \frac{\partial \tilde{\mathbf{h}}_v^{(L)}}{\partial \mathbf{h}_v^{(0)}} \right) \\ \ge \sigma_{\min} \left( \prod_{k=1}^L \mathbf{W}_{\text{skip}}^{(k)} \right) - \|\mathcal{J}_{\text{graph}}(v)\| \\ \ge \prod_{k=1}^L \sigma_{\min}(\mathbf{W}_{\text{skip}}^{(k)}) - \|\mathcal{J}_{\text{graph}}(v)\|\end{gathered}}
\end{equation}
Under the stated coefficient-independence assumptions, absence of every directed return walk to $v$ within the evaluated depth implies $\mathcal{J}_{\text{graph}}(v)=\mathbf{0}$. Then the minimum singular value of the skip product is bounded below by the product of the individual minimum singular values; equality is not generally implied. In cyclic or feature-dependent attention models, the non-skip derivative can reinforce or cancel the skip path. The positive-sensitivity conclusion therefore requires the displayed strict inequality and is not an unconditional guarantee for all GNNs.
\end{proof}

\subsection{Shortcut Behavior and Bilinear Gating}
While Theorem~\ref{thm:theorem1} gives a conditional sensitivity bound for additive residual skips ($\tilde{\mathbf{h}}_v^{(l)} = \mathbf{h}_v^{(l)} + \mathbf{W}_{\text{skip}}^{(l)} \tilde{\mathbf{h}}_v^{(l-1)}$), empirical testing reveals a secondary failure mode: \textit{inductive shortcut collapse}.

\paragraph{The Shortcut Memorization Failure}
Because certificate template attributes (e.g., \texttt{CT\_FLAG\_ENROLLEE\_SUPPLIES\_SUBJECT} $= 1$ and Client Authentication EKU $= 1$) are strongly correlated with positive escalation labels in synthetic training distributions, additive skip connections allow the network to bypass the multi-hop graph aggregation entirely. The MLP classification head learns an inductive shortcut:
\begin{equation}
\fitmath{\hat{y}_t = \text{MLP}(\tilde{\mathbf{h}}_t^{(L)}) \approx \text{MLP} \left( \prod_{k=1}^L \mathbf{W}_{\text{skip}}^{(k)} \mathbf{h}_t^{(0)} \right)}
\end{equation}
When the trained model is evaluated against \textbf{adversarial hard negatives}—templates that possess vulnerable attribute flags but reside in disconnected graph components or require mandatory manager approval—standard additive-skip GNNs memorize the flag pattern and classify them as vulnerable, showing poor agreement on the historical canonical generator. The matched audit in Section~\ref{sec:evaluation} corrects its reachable-path labels and separates augmentation from implementation differences.

\paragraph{Multiplicative Bilinear Reachability Gating (\texttt{CertGraph-Gated}):}
To reduce shortcut memorization, we condition messages on a learned bilinear compatibility gate. The division between learned scoring and symbolic checking follows a broader neuro-symbolic design~\cite{garcez2023neurosymbolic}. Exact reachability is checked separately by the existing symbolic verification step; the learned gate alone is not a Boolean authorization predicate. We formulate the \textbf{Multiplicative Bilinear Reachability Gate}:
\begin{equation}
\fitmath{\mathbf{g}_{uv}^{(l)} = \sigma \left( \mathbf{h}_u^{(l-1)\top} \mathbf{W}_g^{(l)} \mathbf{h}_v^{(l-1)} + b_g^{(l)} \right)}
\end{equation}
where $\mathbf{W}_g^{(l)} \in \mathbb{R}^{d \times d}$ is a bilinear interaction matrix, $b_g^{(l)} \in \mathbb{R}$ is a scalar bias, and $\sigma(\cdot)$ is the sigmoid function mapping to $(0, 1)$.

The gated message vector $\mathbf{m}_{vu}^{(l)}$ and state update are defined as:
\begin{equation}
\fitmath{\begin{gathered}\mathbf{m}_{vu}^{(l)} = \mathbf{g}_{uv}^{(l)} \odot \left( \mathbf{W}_r^{(l)} \mathbf{h}_u^{(l-1)} \right)\end{gathered}}
\end{equation}
\begin{equation}
\fitmath{\begin{gathered}\mathbf{h}_v^{(l)} = \sigma \left( \sum_{r \in \mathcal{R}_{\text{in}}(\phi(v))} \sum_{u \in \mathcal{N}_r(v)} \alpha_{vu}^{(l)} \mathbf{m}_{vu}^{(l)} \right)\\ \quad+\mathbf{g}_{vv}^{(l)} \odot \left( \mathbf{W}_{\text{skip}}^{(l)} \mathbf{h}_v^{(l-1)} \right)\end{gathered}}
\end{equation}
where $\odot$ denotes the element-wise Hadamard product.

\begin{proposition}[Conditional Attenuation by Multiplicative Gating]
For a bounded transformed neighbor message, if its scalar gate tends to zero, multiplicative attenuation gives:
\begin{equation}
\fitmath{\lim_{\mathbf{g}_{uv} \to 0} \mathbf{m}_{vu}^{(l)} = \mathbf{0}}
\end{equation}
Disconnection alone does not imply this gate limit: with zero embeddings and zero bias, the sigmoid gate equals $1/2$. Other residual paths or classifier biases can also produce a positive output. The attenuation statement therefore does not prove universal elimination of shortcuts. The historical numerical implementation uses a feature-channel gate rather than this literal scorer. We now implement a separate two-layer, four-head realization of the displayed bilinear edge and self-skip equations and evaluate it with identical-initialization additive controls and matched augmentation (Table~\ref{tab:literal_bilinear}). All relevant parameter families are trainable through autograd. The two implementations have identical held-out classifications in that new study, and both miss the positive long-path cases. This validates an auditable realization of the formulation, rather than a universal shortcut-elimination or authorization guarantee. A second-tier verifier must process unresolved and out-of-horizon cases, including neural negatives, if it is to recover those observed misses.
\end{proposition}
